\documentclass[11pt]{article}
\usepackage[margin=1.1in]{geometry}
\usepackage{amsmath,amssymb,amsthm,mathtools,booktabs,array,tabularx,longtable}
\usepackage[dvipsnames]{xcolor}
\usepackage[colorlinks=true,linkcolor=MidnightBlue,citecolor=MidnightBlue,urlcolor=MidnightBlue]{hyperref}
\usepackage{fancyhdr}
\pagestyle{fancy}\fancyhf{}\renewcommand{\headrulewidth}{0pt}
\fancyfoot[R]{\thepage}
\fancyfoot[L]{\scriptsize\color{gray}Generated by the Transport Lab from approved parts of its store; not a deposited paper. Papers of record: \href{https://zenodo.org/communities/transport-papers}{Zenodo, Transport Papers}}
\fancypagestyle{plain}{\fancyhf{}\fancyfoot[R]{\thepage}\fancyfoot[L]{\scriptsize\color{gray}Generated by the Transport Lab from approved parts of its store; not a deposited paper. Papers of record: \href{https://zenodo.org/communities/transport-papers}{Zenodo, Transport Papers}}\renewcommand{\headrulewidth}{0pt}}
\newtheorem{theorem}{Theorem}[section]
\newtheorem{proposition}[theorem]{Proposition}
\newtheorem{lemma}[theorem]{Lemma}
\newtheorem{corollary}[theorem]{Corollary}
\newtheorem{conjecture}[theorem]{Conjecture}
\newtheorem{refutation}[theorem]{Refutation}
\newtheorem{prediction}[theorem]{Prediction}
\newtheorem{principle}[theorem]{Principle}
\newtheorem{observation}[theorem]{Observation}
\newtheorem{axiom}[theorem]{Axiom}
\theoremstyle{definition}
\newtheorem{definition}[theorem]{Definition}
\newtheorem{fact}[theorem]{Fact}
\newtheorem{assembled}[theorem]{Assembly}
\newtheorem{postulate}[theorem]{Postulate}
\newtheorem{premise}[theorem]{Premise}
\newtheorem{hypothesis}[theorem]{Hypothesis}
\newtheorem{remark}[theorem]{Remark}
\newtheorem{example}[theorem]{Example}
\theoremstyle{plain}
\setlength{\marginparwidth}{0.8in}
\newcommand{\changedmark}[1]{\marginpar{\raggedright\scriptsize\color{BrickRed}#1}}
\providecommand{\Aa}{}\renewcommand{\Aa}{\mathcal{A}}
\providecommand{\Ad}{}\renewcommand{\Ad}{\widetilde{\mathrm{Ad}}}
\providecommand{\Aev}{}\renewcommand{\Aev}{\mathcal{A}^{+}}
\providecommand{\Aod}{}\renewcommand{\Aod}{\mathcal{A}^{-}}
\providecommand{\B}{}\renewcommand{\B}{\mathrm{B}}
\providecommand{\C}{}\renewcommand{\C}{\mathbb{C}}
\providecommand{\CAR}{}\renewcommand{\CAR}{\mathrm{CAR}}
\providecommand{\CS}{}\renewcommand{\CS}{{\mathbb{S}^{+}}}
\providecommand{\CSm}{}\renewcommand{\CSm}{{\mathbb{S}^{-}}}
\providecommand{\Cl}{}\renewcommand{\Cl}{\mathrm{Cl}}
\providecommand{\Con}{}\renewcommand{\Con}{\mathcal{C}}
\providecommand{\Dl}{}\renewcommand{\Dl}{\Delta}
\providecommand{\EN}{}\renewcommand{\EN}{E_{N}}
\providecommand{\End}{}\renewcommand{\End}{\mathrm{End}}
\providecommand{\F}{}\renewcommand{\F}{\mathbb{F}}
\providecommand{\G}{}\renewcommand{\G}{\Gamma}
\providecommand{\HH}{}\renewcommand{\HH}{\mathbb{H}}
\providecommand{\IS}{}\renewcommand{\IS}{I\sigma_3}
\providecommand{\J}{}\renewcommand{\J}{\mathbf{J}}
\providecommand{\LH}{}\renewcommand{\LH}{L_{H}}
\providecommand{\Mat}{}\renewcommand{\Mat}{\mathrm{M}}
\providecommand{\Mo}{}\renewcommand{\Mo}{M^{\circ}}
\providecommand{\Mp}{}\renewcommand{\Mp}{M^{+}}
\providecommand{\Msun}{}\renewcommand{\Msun}{M_{\odot}}
\providecommand{\Mt}{}\renewcommand{\Mt}{\widetilde{M}}
\providecommand{\Ncone}{}\renewcommand{\Ncone}{\mathcal{N}}
\providecommand{\Neff}{}\renewcommand{\Neff}{N_{\mathrm{eff}}}
\providecommand{\PCS}{}\renewcommand{\PCS}{\mathbb{P}{\mathbb{S}^{+}}}
\providecommand{\PT}{}\renewcommand{\PT}{\mathbb{PT}}
\providecommand{\Pin}{}\renewcommand{\Pin}{\mathrm{Pin}}
\providecommand{\Q}{}\renewcommand{\Q}{\mathbb{Q}}
\providecommand{\R}{}\renewcommand{\R}{\mathbb{R}}
\providecommand{\Ric}{}\renewcommand{\Ric}{\mathrm{Ric}}
\providecommand{\SBH}{}\renewcommand{\SBH}{S_{\mathrm{BH}}}
\providecommand{\Sc}{}\renewcommand{\Sc}{\mathbb{S}}
\providecommand{\Sch}{}\renewcommand{\Sch}{\mathrm{P}}
\providecommand{\Sig}{}\renewcommand{\Sig}{\Sigma}
\providecommand{\Som}{}\renewcommand{\Som}{S_\omega}
\providecommand{\Sp}{}\renewcommand{\Sp}{\mathbb{S}}
\providecommand{\Sph}{}\renewcommand{\Sph}{\mathrm{S}}
\providecommand{\Spi}{}\renewcommand{\Spi}{S_\pi}
\providecommand{\Spin}{}\renewcommand{\Spin}{\mathrm{Spin}}
\providecommand{\Sth}{}\renewcommand{\Sth}{S_{\mathrm{th}}}
\providecommand{\T}{}\renewcommand{\T}{\mathbb{T}}
\providecommand{\TH}{}\renewcommand{\TH}{T_{\mathrm{H}}}
\providecommand{\Tr}{}\renewcommand{\Tr}{\operatorname{Tr}}
\providecommand{\UHF}{}\renewcommand{\UHF}{\mathrm{UHF}}
\providecommand{\Z}{}\renewcommand{\Z}{\mathbb{Z}}
\providecommand{\Zc}{}\renewcommand{\Zc}{\mathrm{Z}}
\providecommand{\cmg}{}\renewcommand{\cmg}{\mathrm{cm^{2}\,g^{-1}}}
\providecommand{\dS}{}\renewcommand{\dS}{\mathrm{dS}}
\providecommand{\dsc}{}\renewcommand{\dsc}{\mathrm{dsc}}
\providecommand{\dtPG}{}\renewcommand{\dtPG}{dt}
\providecommand{\dtS}{}\renewcommand{\dtS}{dt_{\mathrm{S}}}
\providecommand{\ellv}{}\renewcommand{\ellv}{\ell}
\providecommand{\er}{}\renewcommand{\er}{e_r}
\providecommand{\et}{}\renewcommand{\et}{e_t}
\providecommand{\ev}{}\renewcommand{\ev}{\mathrm{ev}}
\providecommand{\gl}{}\renewcommand{\gl}{\mathfrak{gl}}
\providecommand{\hatV}{}\renewcommand{\hatV}{\widehat{V}}
\providecommand{\hb}{}\renewcommand{\hb}{\bar{h}}
\providecommand{\hf}{}\renewcommand{\hf}{\hat{f}}
\providecommand{\id}{}\renewcommand{\id}{\mathrm{id}}
\providecommand{\kpc}{}\renewcommand{\kpc}{\mathrm{kpc}}
\providecommand{\lP}{}\renewcommand{\lP}{\ell_{P}}
\providecommand{\mD}{}\renewcommand{\mD}{\mathbf{m}_{D}}
\providecommand{\mb}{}\renewcommand{\mb}{\bar m}
\providecommand{\mnu}{}\renewcommand{\mnu}{\Sigma m_{\nu}}
\providecommand{\nInf}{}\renewcommand{\nInf}{n_\infty}
\providecommand{\nO}{}\renewcommand{\nO}{n_0}
\providecommand{\ninf}{}\renewcommand{\ninf}{n_{\infty}}
\providecommand{\nz}{}\renewcommand{\nz}{n_{0}}
\providecommand{\proj}{}\renewcommand{\proj}{\operatorname{proj}}
\providecommand{\pv}{}\renewcommand{\pv}{\pi}
\providecommand{\rad}{}\renewcommand{\rad}{\mathrm{rad}}
\providecommand{\sgn}{}\renewcommand{\sgn}{\operatorname{sgn}}
\providecommand{\shA}{}\renewcommand{\shA}{\mathcal{A}}
\providecommand{\shC}{}\renewcommand{\shC}{\mathcal{C}}
\providecommand{\shE}{}\renewcommand{\shE}{\mathcal{E}}
\providecommand{\shJ}{}\renewcommand{\shJ}{\mathcal{J}}
\providecommand{\shK}{}\renewcommand{\shK}{\mathcal{K}}
\providecommand{\shL}{}\renewcommand{\shL}{L}
\providecommand{\shW}{}\renewcommand{\shW}{\mathcal{W}}
\providecommand{\sll}{}\renewcommand{\sll}{\mathfrak{sl}}
\providecommand{\so}{}\renewcommand{\so}{\mathfrak{so}}
\providecommand{\su}{}\renewcommand{\su}{\mathfrak{su}}
\providecommand{\supp}{}\renewcommand{\supp}{\operatorname{supp}}
\providecommand{\tr}{}\renewcommand{\tr}{\operatorname{tr}}
\providecommand{\uu}{}\renewcommand{\uu}{\mathfrak{u}}
\providecommand{\br}{}\renewcommand{\br}[1]{\langle #1\rangle}
\providecommand{\bra}{}\renewcommand{\bra}[1]{\langle #1|}
\providecommand{\file}{}\renewcommand{\file}[1]{\texttt{#1}}
\providecommand{\hp}{}\renewcommand{\hp}[2]{( #1 \,|\, #2 )}
\providecommand{\ip}{}\renewcommand{\ip}[2]{\langle #1, #2\rangle}
\providecommand{\ket}{}\renewcommand{\ket}[1]{|#1\rangle}
\providecommand{\rel}{}\renewcommand{\rel}[2]{D\bigl(#1 \,\|\, #2\bigr)}
\providecommand{\rev}{}\renewcommand{\rev}[1]{\widetilde{#1}}
\providecommand{\slashed}{}\renewcommand{\slashed}[1]{{#1}\mkern-10mu/}
\title{\bfseries Falling In (guide length)}
\author{Leon Butler\\ \normalsize\href{mailto:lbutler.math@gmail.com}{\texttt{lbutler.math@gmail.com}}\\ \normalsize ORCID: \href{https://orcid.org/0009-0006-4449-5617}{0009-0006-4449-5617}}
\date{Generated by the Transport Lab, 27 October 2026}
\begin{document}
\maketitle
\begin{abstract}
\normalsize
At the centre of a non-rotating black hole a curvature measure that no change of description alters grows without bound, so the breakdown is real. Both radial light directions swing inward and merge into one at the centre. The obstruction is a direction with no partner, so the repair enlarges the algebra to supply one. Each step up the resulting tower adds exactly one fermionic mode built from vectors. Inside the hole, Einstein's equations for empty space amount exactly to a marker on the metric staying constant. Coming back down is the same as averaging each observable with its parity-flipped version. For a mode whose mixed state gives both outcomes some probability, the partner making the doubled state pure is the mirror copy of its own algebra. Whatever the starting state, each descended mode of nonzero weight reads as thermal at exactly one temperature. The entropy gained in coming down measures how distinguishable the state is from its descended version, and is never negative. Two premises, on allowed states and on smoothness at the horizon, put every mode at the same Hawking temperature. The algebra alone cannot choose that value: its symmetries only flip the sign of an inverse temperature. At that temperature the Fermi-Dirac filling is the only simple Gaussian state thermal for the boost clock. Given the s-wave modelling convention, its collective boson modes have exactly the Bose-Einstein occupancy, up to an exponentially small correction on a finite window.
\end{abstract}
\newpage
\noindent \textbf{Authorship and disclosure. }The founding idea of the series is the author's: that division by zero in a Clifford algebra should be read not as a prohibition but as a transport into a higher-dimensional algebra, and that the transported division need not be accompanied by an inverse operation. The direction of investigation and the critical review of the results are likewise the author's. The mathematical development and the drafting of the text were carried out in extended collaboration with Claude Fable 5 (Anthropic), an AI system.

\noindent\fbox{\parbox{\dimexpr\linewidth-2\fboxsep-2\fboxrule}{\textbf{This is not a deposited paper.} It was generated by the Transport Lab on 27 October 2026 and assembled from approved parts of its store. It follows results across several papers of the series; each is stated with its proof and credited to the paper that first proved it. The official versions of those papers are on Zenodo, in the Transport Papers community, \url{https://zenodo.org/communities/transport-papers}.}}
\newpage

\section*{Introduction}
This guide follows a black hole from its centre, where the usual mathematics breaks down, to the radiation it gives off. It starts at the singularity, because everything later depends on knowing exactly what has gone wrong there and what kind of repair it calls for. Two short sections then describe the enlarged algebra that makes the repair, called the completion: how much room it has for quantum modes, and what Einstein's equations look like inside it. With the completion in place, the guide follows the return journey back down to the observables of a single mode, and asks what that step keeps and what it throws away. The next section puts a price on the return journey, measured as a gain in entropy. The last section shows how the same ingredients produce a temperature, why the algebra alone cannot set its value, and how the resulting state describes the radiation. Results that the argument uses without proving are gathered at the end and cited rather than shown.

\section{What we assume}

\begin{postulate}[(P): univalence]\label{c:singularity.p-univalence}
The premise (P) requires every physical state $\omega$ of the boundary modes to be even under the grading automorphism $\alpha$, that is $\omega \circ \alpha = \omega$. Equivalently, $\omega$ vanishes on the odd sector $\Aod$, so odd elements of the boundary-mode algebra have zero expectation value.

It is a premise, not a derived result: it restricts which states count as physical, and enters as an input wherever it is invoked.
\end{postulate}

\begin{postulate}[(R): regularity at the fixed locus]\label{c:singularity.r-regularity-fixed-locus}
Premise (R) imposes the KMS condition on each single-mode marginal at a fixed period. For mode $k$, with Killing flow $\alpha_t = \operatorname{Ad}(e^{it\omega_k K_k})$ and $a, b \in \{v_k, v_k^{*}\}$, the entire correlation function $F_{a,b}(t) = \omega(a\,\alpha_t(b))$ is required to continue in imaginary Killing time to a single-valued function on the regular Euclidean section of the Schwarzschild field, its two orderings being the two sides of the cut, which meet after one period: $F_{a,b}(t + ip) = \omega(\alpha_t(b)\,a)$ with $p = 2\pi/\kappa$.

The period is the one that smoothness of the Euclidean section at its tip dictates, a local condition. The condition is imposed on the single-mode marginals only.

On the continuum of the horizon's light rays, for the eternal hole, this content follows from the null ground-state premise (N$_{+}$), which has no temperature in it; the ladder wording here is its discretisation. (R) keeps its status as a premise.
\end{postulate}

\section{The singularity: where division fails}

This section looks at the centre of the black hole, shows that the breakdown there is real, and identifies what kind of breakdown it is.

We begin with a quantity that no change of description can alter.

\begin{proposition}[Non-removability: the invariant argument]\label{c:singularity.non-removability-invariant-argument}
The curvature scalar of the solution is $K := R_{abcd}R^{abcd} = 48M^2/r^6$, assembled from the curvature eigenvalue spectrum (each frame plane contributes four times its eigenvalue squared), and it diverges as $r \to 0$. Under the rotation-gauge law the curvature transforms by conjugation, $\mathcal{R}(B) \mapsto R\,\mathcal{R}(B)\tilde R$, which preserves every scalar contraction, and under position-gauge displacements it transforms by pullback, which relocates the function but cannot remove its unbounded range near the singular locus.

Hence $K$ is invariant under every gauge transformation of the theory, $r$-dependent transformations of arbitrary size included, and no gauge choice removes the degeneration at $r = 0$: the division-by-zero locus is a property of the solution, not of its presentation.
\emph{First proved in} T5, Prop 2.7 (\href{https://doi.org/10.5281/zenodo.22721790}{doi:10.5281/zenodo.22721790}).
\end{proposition}
\noindent\emph{Rests on:} \ref{c:singularity.solution-where-divergence-lies}.

With the breakdown shown to be genuine, we watch what happens to the paths of light on the way in.

\begin{proposition}[The cones fold]\label{c:singularity.cones-fold}
In the Painlev\'e-Gullstrand gauge the metric is $g(a,a) = Q(\hb^{-1}(a))$ with $\hb^{-1} = 1 + \beta P$. For $a = \et + \lambda\,\er$ in the $(t,r)$ plane this gives $\hb^{-1}(a) = \et + (\lambda + \beta)\er$ and $g(a,a) = 1 - (\lambda + \beta)^2$, so the two null directions are $\lambda = -\beta \pm 1$, the radial null curves $dr/dt = -\beta \pm 1$. At the horizon, $\beta = 1$, they are the hovering ray $\lambda = 0$ and an ingoing ray $\lambda = -2$.

As $r \to 0$, $\beta \to \infty$ and both branches, normalised by their $\er$-coefficient, converge to $\er$, tipping inward on the way, while $g(\et,\et) = 1 - \beta^2 \to -\infty$, so no rescaling of the classical field stays finite and nondegenerate. The degeneration is through a null kernel direction in the exact sense of the rescaled-limit result, and it is not removable by gauge choice.
\emph{First proved in} T5, Prop 2.5 (\href{https://doi.org/10.5281/zenodo.22721790}{doi:10.5281/zenodo.22721790}).
\end{proposition}
\noindent\emph{Rests on:} \ref{c:singularity.degeneration-exactly}, \ref{c:singularity.non-removability-invariant-argument}, \ref{c:std-algebra.orthogonal-anticommute}.

Knowing which direction collapses, we can now name the kind of obstruction it is.

\begin{proposition}[Type N, and only type N]\label{c:singularity.type-n-only-type-n}
The obstruction at $r = 0$ is of type N: the degenerating direction $\er$ is null for the limit form and lies in its radical, so it is an isotropic direction without a metric partner, in the degenerate sense that the root taxonomy hands to the sheaf framework, whose localisation theorem supplies the partner by completion on the exceptional set. Type C cannot occur: $\Cl(1,3)$ has $p - q = -2 \equiv 6 \pmod 8$, so it is central simple, $\Cl(1,3) \cong \Mat_2(\HH)$, with no central idempotents for ideal fusion to fuse.

Type P is present, as in every algebra of this size, but demands no transport: its rule is the identity assignment, so it neither licenses nor obstructs the ascent. The operative transport rule is the Witt completion, and no floor of the resulting tower ever develops a type-C obstruction.
\emph{First proved in} T5, Prop 3.1 (\href{https://doi.org/10.5281/zenodo.22721790}{doi:10.5281/zenodo.22721790}).
\end{proposition}
\noindent\emph{Rests on:} \ref{c:sheaf.localised-transport}, \ref{c:sheaf.one-sided-completion}, \ref{c:sheaf.one-sided-completion-singularity}, \ref{c:sheaf.schwarzschild-degeneration-boundary-tame}, \ref{c:singularity.degeneration-exactly}, \ref{c:tower.rigidity-clock-frozen}, \ref{c:transport.existence-classification-type-c}, \ref{c:transport.existence-locus-type-n}, \ref{c:transport.type-p-self-repairing}, \ref{c:primer.classification-real-clifford-algebras}, \ref{c:std-algebra.sta-matrix}.

\section{Quantum capacity of the completion}

This short section asks how much room for quantum modes the completion provides as it is built up.

Since the repair works by enlarging the algebra, the natural next question is how much the enlargement can hold.

\begin{theorem}[Witt index is quantum capacity]\label{c:tower.witt-index-quantum-capacity}
Let $(V, Q)$ be a nondegenerate real quadratic space of signature $(p,q)$, and call a family of vectors $(a_i, b_i)_{i \leq m}$ of $\Cl(V,Q)$ a mode family if it satisfies $a_i a_j + a_j a_i = 0$, $b_i b_j + b_j b_i = 0$ and $a_i b_j + b_j a_i = \delta_{ij}$. The largest possible $m$ is exactly the Witt index $\min(p,q)$. For the upper bound, the relations force the $a_i$ to span a totally isotropic subspace, and pairing against the $b_j$ shows the $a_i$ are linearly independent; a totally isotropic subspace of a nondegenerate form of signature $(p,q)$ has dimension at most $\min(p,q)$. The bound is attained by the normalised null pairs of $\min(p,q)$ orthogonal hyperbolic planes split off by a Witt decomposition.

Along the tower, $T_k$ has signature $(p{+}k, q{+}k)$ and hence Witt index $\min(p,q) + k$, so each rung raises the number of vector-realised modes by exactly one and by no more.
\emph{First proved in} T3, Thm 4.3 (\href{https://doi.org/10.5281/zenodo.22721556}{doi:10.5281/zenodo.22721556}).
\end{theorem}
\noindent\emph{Rests on:} \ref{c:tower.ladder-pairs-satisfy-car-vectors}, \ref{c:std-algebra.vector-squares}, \ref{c:std-algebra.witt-decomposition}.

\section{Einstein's equations in the completion}

Here the question is what form Einstein's equations for empty space take once the completion is in place.

From the algebra of the completion we turn to its geometry.

\begin{theorem}[Covariant form of the vacuum equations]\label{c:interior.covariant-form-vacuum-equations}
Let $g$ be the derived metric of a gravitational field on flat spacetime, on its invertible region $\Mo$, and $I_g = (1,0,0) = \nInf$ its pure-scale tractor. The tractor connection, which on flat spacetime is the $\Cl(2,4)$ gauge field $\Omega_{T}$ built from $\Omega$, the frame $\hb^{-1}$ and the Schouten map, gives $\nabla_a I_g = (0,\, -\Sch_{ab},\, 0)$, so $I_g$ is parallel exactly when $\Sch = 0$, which holds exactly when $g$ is Ricci-flat, which is equivalent to the GTG vacuum field equations. The section is null, $Q(I_g) = 0$.

So the vacuum field equations on $\Mo$ are equivalent to parallelism of the scale tractor attached to $g$. The existence of some parallel null tractor is strictly weaker: it says only that $[g]$ is almost Ricci-flat, and a non-constant rescaling $e^{2\omega}g$ keeps a parallel section while losing the field equations. The statement is covariant under the gauge transformations of $\mathcal{T}$, whose structure algebra $\mathfrak{so}(2,4)$ acts through the bivectors of $\Cl(2,4)$.
\emph{First proved in} T6, Thm 4.5 (\href{https://doi.org/10.5281/zenodo.22721990}{doi:10.5281/zenodo.22721990}).
\end{theorem}
\noindent\emph{Rests on:} \ref{c:interior.imported-equivalence-einstein-equations}, \ref{c:interior.imported-tractor-connection}, \ref{c:interior.metric-parallelism-trace-schouten-condition}.

\section{Descent, parity and states}

Having gone up into the completion, we need a way back down to what a single mode can show us, and this section studies that return step.

With the geometry settled, we look first at the map that comes back down from the completion.

\begin{theorem}[The descent is the parity average]\label{c:singularity.descent-parity-average}
On the single-mode algebra $\Aa$ the descent equals the parity average: $E(x) = \tfrac12\bigl(x + \alpha(x)\bigr) = \tfrac12(x + \sigma x \sigma)$. Hence $E$ is the group average over $\{\id, \alpha\}$ and the unique conditional expectation of $\Aa$ onto its fixed-point algebra $\Aev$.

Its range is $\Aev$, its kernel is exactly the adjoined ladder sector $\Aod$, and it preserves every state with $\omega \circ \alpha = \omega$.
\emph{First proved in} T5, Thm 5.7 (\href{https://doi.org/10.5281/zenodo.22721790}{doi:10.5281/zenodo.22721790}).
\end{theorem}
\noindent\emph{Rests on:} \ref{c:singularity.descent-what-it-erases}.

Because coming down discards information, it is natural to ask where the partner that keeps the larger state pure comes from.

\begin{theorem}[Purity is an output: the partner is the commutant]\label{c:singularity.purity-output-partner-commutant}
Let $\omega$ be a faithful state on the mode factor with $\varrho = \operatorname{diag}(\lambda_0^2, \lambda_1^2)$, $\lambda_i > 0$. Its GNS representation lives on $\C^2 \otimes \C^2$ with cyclic vector $\Omega = \sum_i \lambda_i |ii\rangle$, cyclic and separating for $\pi(\Aa) = \Aa \otimes 1$; the commutant $\pi(\Aa)' = 1 \otimes \Aa = J\pi(\Aa)J$ is the partner copy, and any two purifications of $\omega$ differ by a unitary on the purifying factor alone.

When $\omega$ is the Gibbs state at inverse temperature $\beta$, $\Omega$ is the thermofield double with $\tan\theta = \lambda_1/\lambda_0 = e^{-\omega_k/2T}$, the same vector the Schmidt-forcing argument produces. So purity of the doubled state, the existence of the partner and the thermofield form are consequences of the one-sided state rather than independent choices.
\emph{First proved in} T5, Thm 7.27 (\href{https://doi.org/10.5281/zenodo.22721790}{doi:10.5281/zenodo.22721790}).
\end{theorem}
\noindent\emph{Rests on:} \ref{c:singularity.modular-theory-matrix-factor-proved-outright}.

Turning from the doubled state back to a single mode, we ask what shape its descended state can take.

\begin{theorem}[The descended form is state independent]\label{c:singularity.descended-form-state-independent}
For any state on the boundary mode space and any mode $k$, the descended marginal $E_k(\varrho_k)$ is diagonal in the number basis, because the range of $E_k$ is the even subalgebra. A diagonal density matrix $\operatorname{diag}(p_k, 1 - p_k)$ is then the Gibbs state $e^{-\beta_k\omega_k N_k}/Z_k$ for a unique $\beta_k \in [-\infty, +\infty]$, fixed by $\beta_k\omega_k = \log\bigl(p_k/(1-p_k)\bigr)$, provided $\omega_k \neq 0$.

At a zero-weight mode the right side is $\tfrac12\,\id$ for every $\beta_k$, so uniqueness fails there and existence requires $p_k = \tfrac12$. No hypothesis on the ascended state is used.
\emph{First proved in} T5, Thm 7.6 (\href{https://doi.org/10.5281/zenodo.22721790}{doi:10.5281/zenodo.22721790}).
\end{theorem}
\noindent\emph{Rests on:} \ref{c:singularity.descent-parity-average}.

\section{Entropy and information}

Coming down is not free, and this section measures exactly what it costs.

Knowing what the descent leaves behind, we now put an exact figure on what it removes.

\begin{theorem}[The deficit theorem: the cost of descent, exactly]\label{c:singularity.deficit-theorem-cost-descent-exactly}
Let $E$ be the pinching by a complete family of orthogonal projectors, the multi-mode descent $\EN$ included. Then $E$ is idempotent, trace-preserving and self-adjoint for the Hilbert-Schmidt pairing, and for every density matrix $\rho$ the entropy increase under $E$ is exactly a relative entropy: $S(E(\rho)) - S(\rho) = \rel{\rho}{E(\rho)} \geq 0$, with equality if and only if $\rho = E(\rho)$.

The level discrepancy of the chain is this quantity, so the cost of descent is an exact relative entropy rather than a qualitative statement.
\emph{First proved in} T5, Thm 9.3 (\href{https://doi.org/10.5281/zenodo.22721790}{doi:10.5281/zenodo.22721790}).
\end{theorem}

\section{Temperature and Hawking radiation}

The last section turns the descended states into a temperature, and that temperature into radiation.

With the cost understood, we use the two premises to fix which temperature the descended modes actually have.

\begin{theorem}[The rate, and the discharge of hypothesis (iii)]\label{c:singularity.rate-discharge-hypothesis-iii}
Assume (P) and (R). At each mode the marginal is diagonal in the number basis, $\operatorname{diag}(1 - n, n)$ with $n = \omega(v^*v)$, and the gluing condition of (R) on the ladder pair reduces to $e^{p\omega}n = 1 - n$, so $n = (e^{p\omega} + 1)^{-1}$ with period $p = 2\pi/\kappa$. The marginal is therefore the Gibbs state of $K_k$ at $\beta_k = \beta = 2\pi/\kappa = 8\pi M$, that is $T = \kappa/2\pi = 1/8\pi M = \TH$, with occupancy $\omega(v_k^*v_k) = (e^{\beta\omega_k} + 1)^{-1} \in (0,1)$.

Since the occupancy lies strictly between $0$ and $1$, every marginal is faithful, hence cyclic and separating for its mode factor: the cyclic-and-separating hypothesis of the modular descent theorem is a consequence of (R), not an assumption. The temperature is the same at every mode without any condition on the joint state, because each mode sees the same $\kappa$, which is constant over the bifurcation sphere.
\emph{First proved in} T5, Thm 7.20 (\href{https://doi.org/10.5281/zenodo.22721790}{doi:10.5281/zenodo.22721790}).
\end{theorem}
\noindent\emph{Rests on:} \ref{c:singularity.ascended-marginals-are-diagonal}, \ref{c:singularity.computed}, \ref{c:singularity.descended-form-state-independent}, \ref{c:singularity.modular-structure-completed-pair}, \ref{c:singularity.plane-rate}, \ref{c:singularity.p-univalence}, \ref{c:singularity.r-regularity-fixed-locus}.

Having found the Hawking value, we ask how much of it the algebra itself determines.

\begin{theorem}[Rate rigidity: the temperature is not algebra intrinsic]\label{c:singularity.rate-rigidity-temperature-not-algebra-intrinsic}
Let $\alpha_t = \operatorname{Ad}(e^{itK})$ be the boost's phase flow on the single-mode algebra $\Aa$, the rotor flow $\exp(\lambda\sigma/2)$ continued to imaginary rapidity, with $\beta$-KMS states $\omega_\beta$; every automorphism of this matrix factor is inner, $\theta = \operatorname{Ad}(U)$. An automorphism commuting with the flow fixes every $\omega_\beta$; one reversing the flow sends $\omega_\beta$ to $\omega_{-\beta}$; and a constant reparametrisation $\alpha_t \mapsto \alpha_{ct}$ with $c > 0$ leaves $\omega_\beta$ unchanged while rescaling its inverse temperature to $\beta/c$.

Hence automorphisms intertwining the flow act on the line of KMS states with orbits $\{\beta, -\beta\}$: $|\beta|$ is an invariant of the pair (state, parametrised flow), and the line carries no automorphism-theoretic structure beyond the orientation flip. A condition singling out one point of the line must either stipulate the state or refer to data outside $(\Aa, \alpha)$, and the value of $|\beta|$ is meaningful only relative to a parametrisation of the flow, which the algebra does not supply.
\emph{First proved in} T5, Thm 7.16 (\href{https://doi.org/10.5281/zenodo.22721790}{doi:10.5281/zenodo.22721790}).
\end{theorem}
\noindent\emph{Rests on:} \ref{c:singularity.kms-uniqueness-factor}.

With that caution in mind, we assemble the thermal state across all the boundary modes together.

\begin{theorem}[The I-state]\label{c:hawking.i-state}
Take the boundary-mode algebra $\mathfrak{A}$ with its boost dynamics $\tau$ at $\beta = 1/\TH = 2\pi/\kappa$. The quasi-free state $\omega_I$ with Fermi-Dirac symbol $A_I = (1+e^{\beta h})^{-1}$, that is occupancies $n_j = 1/(e^{\omega_j/\TH}+1)$, restricts on every window $F$ to the Gibbs state of $H_F = \sum_{j\in F}\omega_j N_j$ and is KMS there exactly; it is the unique quasi-free extension of these window states, and it is $\tau$-KMS at $\beta$ on the whole algebra.

Within the quasi-free class the converse holds. Applying the KMS identity to two creation operators shows that the pairing vanishes, so a quasi-free $\tau$-KMS state at $\beta$ is automatically gauge-invariant. Applying it to a creation and an annihilation operator, with the CAR, gives $A(1+e^{-\beta h}) = e^{-\beta h}$, hence $A = A_I$: the two-point data are forced, and $\omega_I$ is the unique quasi-free state whose modular clock is the geometry's clock.
\emph{First proved in} T7, Thm 3.6 (\href{https://doi.org/10.5281/zenodo.22722121}{doi:10.5281/zenodo.22722121}).
\end{theorem}
\noindent\emph{Rests on:} \ref{c:hawking.gibbs-states-are-kms-conversely-finite}, \ref{c:hawking.imported-kms-passes-inductive-limits}, \ref{c:hawking.imported-quasi-free-states-car-algebra}.

It remains to see what this fermionic state looks like when its levels are repackaged as bosons.

\begin{proposition}[Bose-Einstein occupancy, with quantified window corrections]\label{c:hawking.bose-einstein-occupancy-quantified-window-corrections}
On the bilateral ladder with level spacing $\epsilon$ and the sea at level $0$, carrying the linear s-wave weight assignment as a modelling input, the I-state occupancies are $n_k = (x^{-k}+1)^{-1}$ with Boltzmann factor $x = e^{-\beta\epsilon} \in (0,1)$. They obey the exact detailed-balance identity $n_k - n_{k+m} = (x^{-m}-1)\,n_{k+m}(1-n_k)$ for every $k$ and every $m \ge 1$.

Summing over $k \in \Z$, the left side telescopes to $1$ in each of the $m$ residue classes modulo $m$, so $\sum_{k} n_{k+m}(1-n_k) = m/(e^{\beta\epsilon m}-1)$. This sum is the current two-point function $\langle a_m^\dagger a_m\rangle_I$, so the unit-normalised boson $\hat a_m = a_m/\sqrt{m}$ has exactly the Bose-Einstein occupancy $1/(e^{\beta\epsilon m}-1)$ at $\TH$. On the window $[-K, K]$ with the sea at the centre, the bilateral sum exceeds the window sum by a positive amount at most $2m\,x^{K-m}/(1-x)$.
\emph{First proved in} T7, Prop 6.2 (\href{https://doi.org/10.5281/zenodo.22722121}{doi:10.5281/zenodo.22722121}).
\end{proposition}
\noindent\emph{Rests on:} \ref{c:hawking.i-state}, \ref{c:hawking.imported-s-wave-bosonic-sector}.

\section{Results used from elsewhere}
The proofs above use the following results of the series, stated briefly; each is proved in the paper named.

\begin{fact}[Gibbs states are KMS, and conversely, in finite dimensions]\label{c:hawking.gibbs-states-are-kms-conversely-finite}
On a full matrix algebra with dynamics generated by a diagonal Hamiltonian, the Gibbs state at inverse temperature $\beta$ is KMS, and it is the unique KMS state at that $\beta$, degenerate spectrum included.
\emph{Source:} T7, Lem 3.1 (\href{https://doi.org/10.5281/zenodo.22722121}{doi:10.5281/zenodo.22722121}).
\end{fact}

\begin{fact}[Imported: KMS passes to inductive limits]\label{c:hawking.imported-kms-passes-inductive-limits}
A state that is KMS at $\beta$ on each subalgebra of an increasing chain of dynamics-invariant subalgebras is KMS at $\beta$ on the closure of their union.
\emph{Source:} O. Bratteli and D. W. Robinson, Operator Algebras and Quantum Statistical Mechanics 2, 2nd ed., Springer, Berlin (1997)., Prop.~5.3.25.
\end{fact}

\begin{fact}[Imported: quasi-free states of the CAR algebra]\label{c:hawking.imported-quasi-free-states-car-algebra}
Each symbol $0 \le A \le 1$ determines a unique gauge-invariant quasi-free state of the CAR algebra; quasi-free states are fixed by their symbol and pairing, and gauge invariance is the vanishing of the pairing.
\emph{Source:} H. Araki, On quasifree states of CAR and Bogoliubov automorphisms, Publ. RIMS Kyoto Univ. 6 (1971), 385--442.; O. Bratteli and D. W. Robinson, Operator Algebras and Quantum Statistical Mechanics 2, 2nd ed., Springer, Berlin (1997)., \S 5.2.
\end{fact}

\begin{fact}[Imported: the s-wave bosonic sector]\label{c:hawking.imported-s-wave-bosonic-sector}
Boundary-mode currents satisfy the free-boson relations exactly up to an edge term, and under Premise (Q) they carry Bose-Einstein occupancy at $\TH$ in the s-wave sector.
\emph{Source:} T7, Fact 2.4 (\href{https://doi.org/10.5281/zenodo.22722121}{doi:10.5281/zenodo.22722121}).
\end{fact}

\begin{fact}[Imported: equivalence with the Einstein equations]\label{c:interior.imported-equivalence-einstein-equations}
Wherever $\hb$ is invertible, the GTG field equations are equivalent to Einstein's equations for the induced metric, and in vacuum to its Ricci-flatness.
\emph{Source:} A. Lasenby, C. Doran and S. Gull, Gravity, gauge theories and geometric algebra, Phil. Trans. R. Soc. Lond. A 356 (1998), 487--582., \S 4 (\S 4.3, covariant forms of the field equations), App.~C; C. Doran and A. Lasenby, Geometric Algebra for Physicists, Cambridge University Press (2003)., Ch.~13-14.
\end{fact}

\begin{fact}[Imported: the tractor connection]\label{c:interior.imported-tractor-connection}
For a gravitational field on flat spacetime, the tractor bundle of its conformal class (the $\Cl(2,4)$ vectors $\sigma\nInf + \mu + \rho\nO$) carries a canonical metric-preserving connection, given in a chosen scale by an explicit formula involving the Schouten tensor.
\emph{Source:} T. N. Bailey, M. G. Eastwood and A. R. Gover, Thomas's structure bundle for conformal, projective and related structures, Rocky Mountain J. Math. 24 (1994), 1191--1217.; S. N. Curry and A. R. Gover, An introduction to conformal geometry and tractor calculus, with a view towards analysis on manifolds, in Asymptotic Analysis in General Relativity, London Mathematical Society Lecture Note Series 443, Cambridge University Press (2018), 86--170., \S 5.
\end{fact}

\begin{fact}[Metric parallelism, and the trace of the Schouten condition]\label{c:interior.metric-parallelism-trace-schouten-condition}
The tractor gauge field preserves the tractor metric, and a Schouten map proportional to the identity forces the factor to be $R/24$ and the field to be Einstein, so $\Sch = 0$ exactly when $\Ric = 0$.
\emph{Source:} T6, Lem 4.2 (\href{https://doi.org/10.5281/zenodo.22721990}{doi:10.5281/zenodo.22721990}).
\end{fact}

\begin{fact}[Classification of the real Clifford algebras]\label{c:primer.classification-real-clifford-algebras}
$\Cl(p,q)$ is a matrix algebra over $\R$, $\C$ or $\HH$, or a sum of two such, according to $p - q \bmod 8$, and it is simple except when $p - q \equiv 1 \pmod 4$.
\emph{Source:} T0, Thm 9.3 (\href{https://doi.org/10.5281/zenodo.22714501}{doi:10.5281/zenodo.22714501}).
\end{fact}

\begin{fact}[Localised transport]\label{c:sheaf.localised-transport}
The degenerate Clifford sheaf embeds canonically in the completed sheaf, and on the wall the radical covector acquires a partner with which it forms a ladder pair that exists exactly over the wall.
\emph{Source:} T4, Cor 4.5 (\href{https://doi.org/10.5281/zenodo.22721669}{doi:10.5281/zenodo.22721669}).
\end{fact}

\begin{fact}[One-sided completion]\label{c:sheaf.one-sided-completion}
Under tameness, with signature $(p,q)$, $p \geq 1$, on the physical side and a trivial radical line bundle, coupling an adjoined hyperbolic plane to the dying direction gives an ambient form of constant signature $(p{+}1, q{+}1)$ on the whole physical side, wall included.
\emph{Source:} T4, Thm 4.4 (\href{https://doi.org/10.5281/zenodo.22721669}{doi:10.5281/zenodo.22721669}).
\end{fact}

\begin{fact}[One-sided completion at the singularity]\label{c:sheaf.one-sided-completion-singularity}
At the blown-up Schwarzschild singularity, coupling an adjoined hyperbolic plane to the dying direction with $\mu\nu < 0$ and bump inside the horizon gives an ambient form of constant signature $(2,4)$ everywhere, boundary included, with no shadow wall and a ladder pair existing exactly over the boundary.
\emph{Source:} T4, Thm 8.12 (\href{https://doi.org/10.5281/zenodo.22721669}{doi:10.5281/zenodo.22721669}).
\end{fact}

\begin{fact}[The Schwarzschild degeneration is boundary-tame in $\shW$]\label{c:sheaf.schwarzschild-degeneration-boundary-tame}
In the weighted module at the blown-up Schwarzschild singularity, the derived cometric of the Painlev\'e-Gullstrand gauge field extends smoothly to the boundary and degenerates there with corank exactly one and to exactly first order, with radical $k = \rho + m\,\dtPG$ annihilating the infalling congruence.
\emph{Source:} T4, Thm 8.7 (\href{https://doi.org/10.5281/zenodo.22721669}{doi:10.5281/zenodo.22721669}).
\end{fact}

\begin{fact}[Ascended marginals are diagonal]\label{c:singularity.ascended-marginals-are-diagonal}
Under univalence every single-mode marginal is diagonal in the number basis, hence equal to its descent and, at nonzero weight, a Gibbs state of the boost generator at a unique $\beta_k$.
\emph{Source:} T5, Thm 7.10 (\href{https://doi.org/10.5281/zenodo.22721790}{doi:10.5281/zenodo.22721790}).
\end{fact}

\begin{fact}[$\kappa = 1/4M$, computed]\label{c:singularity.computed}
For the Schwarzschild field on flat spacetime the redshifted acceleration of static observers, $v = K/\sqrt f$, is $f'(r)/2$, so the surface gravity is $\kappa = 1/4M$.
\emph{Source:} T5, Prop A.2 (\href{https://doi.org/10.5281/zenodo.22721790}{doi:10.5281/zenodo.22721790}).
\end{fact}

\begin{fact}[The degeneration, exactly]\label{c:singularity.degeneration-exactly}
As $r \to 0$ the rescaled gauge map and metric converge exactly to rank-one limits whose common kernel $\et^{\perp}$ contains the folded direction $\er$, null for the limit form.
\emph{Source:} T5, Prop 2.6 (\href{https://doi.org/10.5281/zenodo.22721790}{doi:10.5281/zenodo.22721790}).
\end{fact}

\begin{fact}[The descent, and what it erases]\label{c:singularity.descent-what-it-erases}
The descent $E(x) = PxP + (1-P)x(1-P)$, with $P = v^* v$, keeps a mode's diagonal occupation data and erases its coherences.
\emph{Source:} T5, Prop 5.5 (\href{https://doi.org/10.5281/zenodo.22721790}{doi:10.5281/zenodo.22721790}).
\end{fact}

\begin{fact}[KMS uniqueness at a factor]\label{c:singularity.kms-uniqueness-factor}
On a matrix factor, for each nonzero inverse temperature the only KMS state of an inner flow is the Gibbs state of its generator.
\emph{Source:} T5, Thm 7.13 (\href{https://doi.org/10.5281/zenodo.22721790}{doi:10.5281/zenodo.22721790}).
\end{fact}

\begin{fact}[Modular structure of the completed pair]\label{c:singularity.modular-structure-completed-pair}
The element $\sigma = v^*v - vv^*$ of the completed pair is at once the boost generator, the shifted number operator and the generator of the descent projectors, and its boost flow is the only real rotor flow of the atom.
\emph{Source:} T5, Thm 6.3 (\href{https://doi.org/10.5281/zenodo.22721790}{doi:10.5281/zenodo.22721790}).
\end{fact}

\begin{fact}[Modular theory of a matrix factor, proved outright]\label{c:singularity.modular-theory-matrix-factor-proved-outright}
For a full matrix algebra with a cyclic separating vector, the modular operator, modular conjugation and modular flow are computed explicitly, and the state is KMS for that flow.
\emph{Source:} T5, Lem 7.1 (\href{https://doi.org/10.5281/zenodo.22721790}{doi:10.5281/zenodo.22721790}).
\end{fact}

\begin{fact}[The plane and the rate]\label{c:singularity.plane-rate}
The curvature selects the boost plane of the completed pair invariantly, and at the horizon its rate is the surface gravity $\kappa = 1/4M$, so $2\pi$ of rapidity corresponds to $8\pi M$ of Schwarzschild time.
\emph{Source:} T5, Prop 6.2 (\href{https://doi.org/10.5281/zenodo.22721790}{doi:10.5281/zenodo.22721790}).
\end{fact}

\begin{fact}[The solution, and where the divergence lies]\label{c:singularity.solution-where-divergence-lies}
The gauge-theory Schwarzschild solution satisfies the vacuum equations at every $r > 0$ with $\det \hb = 1$ throughout, so its divergence lies in the connection and the curvature, not in $\hb$.
\emph{Source:} T5, Prop 2.4 (\href{https://doi.org/10.5281/zenodo.22721790}{doi:10.5281/zenodo.22721790}).
\end{fact}

\begin{fact}[Orthogonal vectors anticommute]\label{c:std-algebra.orthogonal-anticommute}
$uv = -vu$ iff $B(u,v) = 0$; in an orthogonal basis $e_I e_J = \pm \prod_{I \cap J} Q(e_i)\; e_{I \triangle J}$ with the sorting sign.
\emph{Source:} the algebra core of the Transport Lab, proved there with an exact check; standard, see C. Doran and A. Lasenby, Geometric Algebra for Physicists, Cambridge University Press (2003), Ch. 4 (foundations of geometric algebra).
\end{fact}

\begin{fact}[The spacetime algebra as a matrix algebra, and the spacetime split]\label{c:std-algebra.sta-matrix}
$\Cl(1,3) \cong \Mat_2(\HH)$, the Dirac matrices represent $\C \otimes \Cl(1,3) \cong \Mat_4(\C)$, and $\sigma_k = \gamma_k\gamma_0$ give $\Cl^{+}(1,3) \cong \Cl(3,0)$ with $\sigma_1\sigma_2\sigma_3 = I$.
\emph{Source:} the algebra core of the Transport Lab, proved there with an exact check; standard, see C. Doran and A. Lasenby, Geometric Algebra for Physicists, Cambridge University Press (2003), Ch. 5 (relativity and spacetime).
\end{fact}

\begin{fact}[The product of two vectors]\label{c:std-algebra.vector-squares}
$v^{2} = Q(v)$ and $uv + vu = 2B(u, v)$, so $uv = B(u, v) + \tfrac12(uv - vu)$.
\emph{Source:} the algebra core of the Transport Lab, proved there with an exact check; standard, see C. Doran and A. Lasenby, Geometric Algebra for Physicists, Cambridge University Press (2003), Ch. 4 (foundations of geometric algebra).
\end{fact}

\begin{fact}[Hyperbolic pairs and the Witt decomposition]\label{c:std-algebra.witt-decomposition}
A null vector in a nondegenerate plane has a unique hyperbolic partner, and $V = H_1 \perp \dots \perp H_{\min(p,q)} \perp A \perp \operatorname{rad} V$ with $A$ definite; maximal totally null subspaces have dimension $\min(p,q) + r$.
\emph{Source:} the algebra core of the Transport Lab, proved there with an exact check; standard, see T. Y. Lam, Introduction to Quadratic Forms over Fields, Graduate Studies in Mathematics 67, American Mathematical Society (2005), Ch. I (hyperbolic planes, the Witt decomposition, Witt's theorems).
\end{fact}

\begin{fact}[The ladder pairs satisfy the CAR, as vectors]\label{c:tower.ladder-pairs-satisfy-car-vectors}
In the $k$-th tower algebra the $k$ adjoined null pairs satisfy the canonical anticommutation relations of $k$ fermionic modes exactly, as vectors, and each new rung adds one mode without altering the existing relations.
\emph{Source:} T3, Prop 4.1 (\href{https://doi.org/10.5281/zenodo.22721556}{doi:10.5281/zenodo.22721556}).
\end{fact}

\begin{fact}[Rigidity; the clock is frozen]\label{c:tower.rigidity-clock-frozen}
Every floor of the tower shares the ground floor's charge, division-ring type, centre, simplicity class, Morita class and obstruction taxonomy, so the physical tower $\Cl(1{+}k, 3{+}k) \cong \Mat_{2^{k+1}}(\HH)$ is frozen at clock position $6$.
\emph{Source:} T3, Thm 3.6 (\href{https://doi.org/10.5281/zenodo.22721556}{doi:10.5281/zenodo.22721556}).
\end{fact}

\begin{fact}[Existence and classification of type C]\label{c:transport.existence-classification-type-c}
$\Cl(p,q)$, $n \geq 1$, has a nontrivial central idempotent exactly when $p - q \equiv 1 \pmod 4$, and then the only ones are $\varepsilon_\pm = \tfrac12(1 \pm \omega)$, which split the algebra into two simple ideals and are exchanged by the grade involution but by no inner automorphism.
\emph{Source:} T1, Thm 3.6 (\href{https://doi.org/10.5281/zenodo.21743863}{doi:10.5281/zenodo.21743863}).
\end{fact}

\begin{fact}[Existence locus of type N]\label{c:transport.existence-locus-type-n}
$\Cl(p,q)$ contains null vectors exactly when $pq \neq 0$; they then form a cone of dimension $n - 1$, and each is a two-sided zero divisor with no left or right inverse in any unital algebra containing $\Cl(p,q)$.
\emph{Source:} T1, Prop 3.3 (\href{https://doi.org/10.5281/zenodo.21743863}{doi:10.5281/zenodo.21743863}).
\end{fact}

\begin{fact}[Type P is self-repairing]\label{c:transport.type-p-self-repairing}
In any unital ring the identity assignment $\sigma(e) = e$ on idempotents satisfies both von Neumann identities and is equivariant under every automorphism of every extension preserving the ring, so type P obstructions are repaired at signature increment $(0,0)$.
\emph{Source:} T1, Prop 6.1 (\href{https://doi.org/10.5281/zenodo.21743863}{doi:10.5281/zenodo.21743863}).
\end{fact}

\begin{prediction}[The no-echo prediction, registered]\label{c:scrambling.no-echo-prediction-registered}
Under Postulate (C1), which puts the quantum structure at the stretched singularity, post-merger ringdown is the unmodified Kerr quasinormal spectrum, with no echo train at the delay $8M\ln(M/m_P)$.
\emph{Kill condition:} One statistically confirmed echo train at the Planck-structure delay.
\end{prediction}

\section*{Conclusion}
The breakdown at the centre of a non-rotating black hole is a feature of the solution itself, not of how it is written down. It is the kind of obstruction caused by a direction with no partner, so the repair is to enlarge the algebra and supply one. Coming back down from the enlargement costs exactly an entropy gain, which is never negative and is zero only when there was nothing to discard. Under the two premises every mode comes out at the same Hawking temperature, with no assumption about correlations between modes. The Fermi-Dirac state at that temperature is the only simple Gaussian state that is thermal for the boost clock of the geometry. Given the s-wave modelling convention, its collective boson modes on the infinite ladder have exactly the Bose-Einstein thermal occupancy. What we assume comes down to a short list. The first premise, univalence, says that physical states of the boundary modes give zero average to every quantity that changes fermion number by an odd amount. The second, regularity, asks each boundary mode on its own to look thermal for the black hole's time flow, at the temperature for which the imaginary-time version of the spacetime is smooth. The s-wave convention behind the boson reading is a modelling input rather than a derived result. The value of the temperature also has meaning only once the clock's parameter is converted into physical time, and the algebra does not supply that conversion. One consequence faces observation directly: merging black holes should ring exactly as general relativity says, with no echoes, and a single confirmed echo would refute the framework.

\end{document}